\subsection{Fehler}
Es wäre sehr hilfreich in der Anleitung zu integrieren, dass man jedes Mal die Div-Werte aufschreiben muss, um im Nachhinein die Fehler ausrechnen zu können. Bei den meisten Signalen ging es noch über die Bilder, jedoch sind nicht immer die Einstellungen des Bildes auch die gewesen, unter denen die Cursor bewegt wurden.

\subsection{Signal 9 - Fouriertransformation}
Die Werte im Messprotokoll sind kompletter Schwachsinn, weshalb ich sie erneut in der Grafik abgelesen habe. Die hohe Sigma-Abweichung für \(f_1\) muss dennoch irgendwie mit dem Ablesen zusammenhängen, obwohl ich es mehrmals gecheckt habe. 

\subsection{Pulsweitenmodulation}
Wahrscheinlich war die Kurve nicht richtig auf den Referenzpunkt der Unterkante der Pulse gesetzt, sondern irgendwie in der Mitte, weswegen der Mittelwert fast \qty{0}{\volt} beträgt. Das kann ich mangels Bildes leider nicht überprüfen. Daran liegt sehr wahrscheinlich die komplett nichtsaussagenden und von einem anderen Planeten kommenden Werte vom Oszi (I mean das Oszi ist so krass, dass kommt auch von nem anderen Planeten. Die Schuld liegt da bei mir). 

Vielleicht wäre es hilfreich am Anfang kurz die verschiedenen Kopplungen zu erklären und wann man den Referenzpunkt richtig bzw. neu setzen muss.

\subsection{Fazit}
In Zukunft könnte man es sich sparen, alle Werte per Hand auszulesen, und mit den automatischen Anzeigen des Oszilloskops arbeiten, und diese evtl. stichprobenartig überprüfen.\\
Das Lernen der Funktionen und die richtige Bedienung entsprechend der Aufgaben zu lernen, hat gedauert. Deswegen war der Versuch bisschen stressig. Die Erkenntnis, dass wir nie Div-Werte aufgeschrieben hat dem den Rest gegeben.\\
Der Versuch ist eigentlich richtig cool, weil das Oszi ein so cooles Gerät ist, aber am Ende war ich fertig. Dafür waren am Ende die Funktionen/Bedienung so klar, dass es dann sehr schnell ging, paar Aufgaben nochmal zu wiederholen. Ich hab aber im Messprotokoll nicht deutlich gemacht, welche Fehlerwerte erfunden und welche nochmal nachgeholt worden sind, weswegen ich soviel als möglich neu über die Bilder erhoben habe.
